\documentclass[11pt,a4paper]{article}

\usepackage[margin=1in]{geometry}
\usepackage{amsmath}
\usepackage{amssymb}
\usepackage{booktabs}
\usepackage{hyperref}
\usepackage{xcolor}
\usepackage{enumitem}
\usepackage{microtype}

\hypersetup{
  colorlinks=true,
  linkcolor=blue,
  citecolor=blue,
  urlcolor=blue
}

\title{Text-Conditioned Music Structure Analysis with Continuous Label Embeddings}
\author{TISMIR Structure Project}
\date{\today}

\newcommand{\audioseq}{\mathbf{A}}
\newcommand{\textset}{\mathbf{E}}
\newcommand{\audiohid}{\mathbf{X}}
\newcommand{\texthid}{\mathbf{Z}}
\newcommand{\logits}{\mathbf{S}}
\newcommand{\probs}{\mathbf{P}}
\newcommand{\labels}{\mathcal{L}}
\newcommand{\frames}{\mathcal{T}}
\newcommand{\ignore}{\mathtt{ignore}}

\begin{document}
\maketitle

\begin{abstract}
This document describes the current model used in the TISMIR structure-analysis project.
The goal is beat-synchronous music structure segmentation conditioned on an arbitrary
set of text labels. Instead of predicting labels from a fixed taxonomy, the model maps
audio frames and text label embeddings into a shared space and predicts frame-level
label assignments from their similarities. The current implementation is inspired by
FACT, the Frame-Action Cross-Attention Temporal Modeling architecture for action
segmentation, but adapts it to open-vocabulary music section labels. The document also
summarizes the optional FACT-inspired modules recently added to the codebase:
intermediate supervision, cross-attention alignment supervision, and attention-fused
prediction.
\end{abstract}

\section{Problem Setting}

Let a song be represented by a beat-synchronous sequence of precomputed audio
embeddings
\begin{equation}
  \audioseq = (a_1,\ldots,a_T), \qquad a_t \in \mathbb{R}^{D_a},
\end{equation}
where $T$ is the number of beat intervals. For the same track, a candidate label set is
provided:
\begin{equation}
  \labels = \{\ell_1,\ldots,\ell_K\}.
\end{equation}
Each label is encoded by an off-the-shelf text encoder, giving
\begin{equation}
  \textset = (e_1,\ldots,e_K), \qquad e_k \in \mathbb{R}^{D_t}.
\end{equation}

The model predicts a distribution over the provided candidate labels for each beat:
\begin{equation}
  p_\theta(y_t = k \mid \audioseq, \labels), \qquad t \in \{1,\ldots,T\}, \quad
  k \in \{1,\ldots,K\}.
\end{equation}
The key design constraint is that the label set is not fixed at training or inference
time. At inference, the user may provide a coarse label set such as
\begin{quote}
\texttt{intro, verse, chorus, outro}
\end{quote}
or a more detailed label set such as
\begin{quote}
\texttt{intro, verse 1, chorus 1, verse 2, chorus 2, bridge, chorus 3, outro}.
\end{quote}
The resulting segmentation should follow the granularity of the provided label set.

\section{Input Preprocessing}

\subsection{Audio Embeddings}

Audio is processed offline. A music audio encoder such as MERT is applied to the full
track, and embeddings are pooled over beat intervals. The default pooling is mean
pooling over each beat interval. The resulting matrix is saved as a NumPy array:
\begin{equation}
  \audioseq \in \mathbb{R}^{T \times D_a}.
\end{equation}
Beat positions are estimated by an external beat tracker and stored alongside the
embeddings.

\subsection{Text Embeddings}

Each section label is normalized, optionally prompted, then encoded by a pretrained
text encoder. The project currently supports several prompt modes. The most relevant
for occurrence-level labels is:
\begin{quote}
\texttt{Music structure label: \{label\}. Meaning: \{occurrence\_prompt\}. Use this label for frames belonging to \{occurrence\_prompt\}.}
\end{quote}
For example:
\begin{quote}
\texttt{Music structure label: chorus 2. Meaning: the second occurrence of the chorus section in this song. Use this label for frames belonging to the second occurrence of the chorus section in this song.}
\end{quote}

\subsection{Annotation Policies}

The codebase supports multiple annotation-processing policies:
\begin{itemize}[leftmargin=*]
  \item \texttt{keep}: use labels as they appear in the original annotation.
  \item \texttt{merge}: merge consecutive identical section labels.
  \item \texttt{enumerate\_all\_occurrences}: enumerate repeated labels in chronological order.
  \item \texttt{enumerate\_base\_occurrences}: strip existing occurrence markers and re-enumerate base labels.
  \item \texttt{enumerate\_consecutive\_repeats}: enumerate only labels involved in consecutive repeated sections.
\end{itemize}

Training may also use a random annotation policy. In that setting, each track and epoch
uses one sampled policy, while validation remains deterministic. This acts as label
granularity augmentation.

\section{Base Architecture}

\subsection{Projection}

Audio and text embeddings are first projected to a common model dimension $d$:
\begin{align}
  \audiohid^{(0)} &= f_a(\audioseq) \in \mathbb{R}^{T \times d}, \\
  \texthid^{(0)} &= f_t(\textset) \in \mathbb{R}^{K \times d}.
\end{align}
The projections are small feed-forward networks, typically linear--GELU--linear.

\subsection{Audio Positional Encoding}

The audio side may use absolute sinusoidal positional encoding:
\begin{equation}
  \audiohid^{(0)} \leftarrow \audiohid^{(0)} + \mathrm{PE}.
\end{equation}
Rotary positional encoding is also implemented, but in recent experiments sinusoidal
encoding has generalized better.

\subsection{Initial Audio and Text Encoders}

The audio sequence is processed by a Transformer encoder:
\begin{equation}
  \audiohid = \mathrm{AudioEnc}(\audiohid^{(0)}).
\end{equation}
For non-bidirectional cross-attention models, the text tokens are processed by a text
Transformer. For the FACT-style bidirectional model, text tokens are first projected,
then updated by an input block using audio information.

\section{FACT-Inspired Bidirectional Blocks}

The model is inspired by FACT~\cite{lu2024fact}, which alternates between frame tokens
and task-level tokens. In this project, the task-level tokens are section label tokens.
The main analogy is:
\begin{center}
\begin{tabular}{lll}
\toprule
FACT terminology & Music-structure model & Meaning \\
\midrule
Frame tokens & Audio beat tokens & Beat-synchronous audio sequence \\
Task-level tokens & Section label tokens & Candidate structure labels \\
Temporal segmentation & Section segmentation & Frame-level label assignment \\
\bottomrule
\end{tabular}
\end{center}

\subsection{Input Block}

The input block lets text label tokens attend to the audio sequence:
\begin{equation}
  \texthid \leftarrow
  \mathrm{LN}\left(
    \texthid + \mathrm{MHA}(\texthid, \audiohid, \audiohid)
  \right),
\end{equation}
then applies a section-token branch Transformer to refine the label tokens.

\subsection{Update Blocks}

Each update block has two cross-attention directions:
\begin{align}
  \texthid &\leftarrow
  \mathrm{SectionBranch}\left(
    \mathrm{LN}\left(
      \texthid + \mathrm{MHA}(\texthid, \audiohid, \audiohid)
    \right)
  \right), \\
  \audiohid &\leftarrow
  \mathrm{FrameBranch}\left(
    \mathrm{LN}\left(
      \audiohid + \mathrm{MHA}(\audiohid, \texthid, \texthid)
    \right)
  \right).
\end{align}
Here $\mathrm{MHA}(Q,K,V)$ denotes multi-head attention. In the current
implementation, the frame branch is a self-attention Transformer encoder rather than a
temporal convolution.

\section{Prediction Head}

Final audio and text tokens are optionally normalized:
\begin{align}
  \hat{x}_t &= \frac{x_t}{\|x_t\|_2}, \\
  \hat{z}_k &= \frac{z_k}{\|z_k\|_2}.
\end{align}
The direct similarity logits are:
\begin{equation}
  s_{t,k}^{\mathrm{direct}}
  =
  \frac{\hat{x}_t^\top \hat{z}_k}{\tau},
\end{equation}
where $\tau$ is a temperature.

The frame-label probability is:
\begin{equation}
  p_t(k) =
  \mathrm{softmax}_k(s_{t,k}).
\end{equation}

\section{Optional FACT-Inspired Improvements}

Three optional modules are now controlled by config flags.

\subsection{Intermediate Supervision}

Each FACT update block can emit intermediate logits:
\begin{equation}
  \logits^{(b)} \in \mathbb{R}^{T \times K}, \qquad b = 1,\ldots,B.
\end{equation}
An auxiliary frame-label loss can be applied at every update block:
\begin{equation}
  \mathcal{L}_{\mathrm{intermediate}}
  =
  \frac{1}{B}\sum_{b=1}^{B}
  \mathrm{CE}(\logits^{(b)}, y).
\end{equation}
This encourages every update stage to produce useful frame-label predictions, rather
than relying only on the final block.

Config:
\begin{verbatim}
model:
  cross_attention:
    intermediate_logits: true

loss:
  intermediate_frame_label:
    weight: 0.2
\end{verbatim}

\subsection{Cross-Attention Alignment Supervision}

The model can expose two attention maps per update block:
\begin{align}
  A^{\mathrm{frame}\rightarrow\mathrm{text}} &\in \mathbb{R}^{T \times K}, \\
  A^{\mathrm{text}\rightarrow\mathrm{frame}} &\in \mathbb{R}^{K \times T}.
\end{align}

Frame-to-text alignment asks each frame to attend to its ground-truth label token:
\begin{equation}
  \mathcal{L}_{f \rightarrow l}
  =
  -\frac{1}{|\frames|}
  \sum_{t \in \frames}
  \log A^{\mathrm{frame}\rightarrow\mathrm{text}}_{t,y_t}.
\end{equation}

Text-to-frame alignment asks each label token to attend to frames assigned to that
label:
\begin{equation}
  \mathcal{L}_{l \rightarrow f}
  =
  -\frac{1}{|\labels|}
  \sum_{k \in \labels}
  \frac{1}{|\frames_k|}
  \sum_{t \in \frames_k}
  \log A^{\mathrm{text}\rightarrow\mathrm{frame}}_{k,t}.
\end{equation}

The combined alignment loss is:
\begin{equation}
  \mathcal{L}_{\mathrm{attn}}
  =
  \lambda_f \mathcal{L}_{f \rightarrow l}
  +
  \lambda_l \mathcal{L}_{l \rightarrow f}.
\end{equation}

Config:
\begin{verbatim}
model:
  cross_attention:
    return_attention: true

loss:
  cross_attention_alignment:
    weight: 0.05
    frame_to_text_weight: 1.0
    label_to_frame_weight: 1.0
    use_all_blocks: true
\end{verbatim}

\subsection{Attention-Fused Prediction}

The frame-to-text attention map can be used as a second prediction path. If
$A_{t,k}$ is the final frame-to-text attention probability, define:
\begin{equation}
  s_{t,k}^{\mathrm{attn}} = \log(A_{t,k} + \epsilon).
\end{equation}
The final logits are:
\begin{equation}
  s_{t,k}
  =
  (1-\alpha)s_{t,k}^{\mathrm{direct}}
  +
  \alpha s_{t,k}^{\mathrm{attn}}.
\end{equation}
This gives the model a direct way to use cross-attention alignment during prediction.

Config:
\begin{verbatim}
model:
  cross_attention:
    attention_fusion:
      enabled: true
      weight: 0.25
\end{verbatim}

\section{Training Objective}

The default supervised objective is frame-label cross-entropy:
\begin{equation}
  \mathcal{L}_{\mathrm{CE}}
  =
  - \sum_{t : y_t \ne \ignore}
  \log p_t(y_t).
\end{equation}

The pairwise probability loss operates on the predicted frame-label distributions:
\begin{equation}
  Q = \probs \probs^\top,
\end{equation}
where
\begin{equation}
  Q_{i,j} = \sum_{k=1}^{K} p_i(k)p_j(k)
\end{equation}
is the model-implied probability that frames $i$ and $j$ share the same label. The
reference matrix is:
\begin{equation}
  R_{i,j} = \mathbb{1}[y_i = y_j].
\end{equation}
The loss is binary cross-entropy between $Q$ and $R$, optionally balanced by class
counts.

The total objective may combine:
\begin{equation}
  \mathcal{L}
  =
  \lambda_{\mathrm{CE}}\mathcal{L}_{\mathrm{CE}}
  +
  \lambda_{\mathrm{pair}}\mathcal{L}_{\mathrm{pair}}
  +
  \lambda_{\mathrm{inter}}\mathcal{L}_{\mathrm{intermediate}}
  +
  \lambda_{\mathrm{attn}}\mathcal{L}_{\mathrm{attn}}.
\end{equation}

\section{Inference}

At inference time, a track is processed with a chosen candidate label set. The model
produces beat-level logits, which may be smoothed temporally:
\begin{equation}
  \tilde{\logits} = \mathrm{Smooth}(\logits).
\end{equation}
Labels are decoded either by simple argmax or by Viterbi decoding with a transition
penalty. Consecutive frames assigned to the same label are merged into segments.

\section{Current Observations}

Recent diagnostics suggest that the model can learn meaningful section structure, but
generalization remains the main challenge. Easy tracks show clean block-diagonal
probability self-similarity matrices. Harder tracks tend to show weak audio-text
margins and label competition in ambiguous middle or late sections. Random annotation
policy training improves full-validation performance by exposing the model to multiple
granularity conventions.

The most promising next directions are:
\begin{enumerate}[leftmargin=*]
  \item supervised cross-attention alignment;
  \item intermediate block-level supervision;
  \item attention-fused prediction;
  \item memory-efficient frame branches for long songs;
  \item lightweight audio-feature augmentation.
\end{enumerate}

\section{Example Experiment Configuration}

\begin{verbatim}
model:
  name: temporal_text_adapter
  audio:
    hidden_dim: 128
    positional_encoding:
      type: sinusoidal
    num_layers: 2
  text:
    hidden_dim: 128
    num_layers: 1
  adapter:
    model_dim: 128
    num_heads: 8
    feedforward_dim: 512
    dropout: 0.2
  cross_attention:
    enabled: true
    bidirectional: true
    num_layers: 4
    num_heads: 8
    intermediate_logits: true
    return_attention: true
    attention_fusion:
      enabled: true
      weight: 0.25

loss:
  frame_label_weight: 1.0
  intermediate_frame_label:
    weight: 0.2
  cross_attention_alignment:
    weight: 0.05
    frame_to_text_weight: 1.0
    label_to_frame_weight: 1.0
    use_all_blocks: true
  pairwise_probability:
    weight: 0.05
    balance: true
\end{verbatim}

\begin{thebibliography}{9}

\bibitem{lu2024fact}
Lu, Y., Liu, J., Zhang, Y., Liu, Y., and Chen, X.
\newblock FACT: Frame-Action Cross-Attention Temporal Modeling for Efficient Action Segmentation.
\newblock In \emph{Proceedings of the IEEE/CVF Conference on Computer Vision and Pattern Recognition}, 2024.
\newblock \url{https://openaccess.thecvf.com/content/CVPR2024/papers/Lu_FACT_Frame-Action_Cross-Attention_Temporal_Modeling_for_Efficient_Action_Segmentation_CVPR_2024_paper.pdf}

\end{thebibliography}

\end{document}
